\section{Personalisierung des Verhältniswahlrechts}

\subsection{Keine parteilosen Kandidaten}

Der Grundgedanke hinter dem Verhältniswahlrecht ist, dass möglichst alle gesellschaftlich
relevanten Geistesströmungen und sozialen Gruppen im Parlament repräsentiert sein sollen. Die
Repräsentation findet dabei durch Parteien statt, die ursprünglich als Gruppen Gleichgesinnter
bezüglich bestimmter Ideen gedacht waren. Es kann allerdings auch vorkommen, dass zwei Parteien
sich inhaltlich kaum unterscheiden, sondern einfach nur zwei Mannschaften sind, die versprechen,
dieselbe Sache besser zu machen. Sie sind zwei Telekommunikationsunternehmen vergleichbar, die
um die Gunst der Kunden wetteifern, indem sie versuchen, in derselben Sache, Telekommunikation,
besser zu sein als die Konkurrenz. Darin ist prinzipiell nichts Schlechtes, solange trotzdem
sichergestellt ist, dass alle relevanten gesellschaftlichen Strömungen ggf.\ durch andere Parteien
noch im Parlament vertreten sind.

Vom Prinzip her könnte auch eine Einzelperson eine ganze Geistesströmung repräsentieren.
Allerdings würde ich aus zwei Gründen trotzdem keine Einzelpersonen zu Länder- und
Bundestagswahlen kandidieren lassen:
\begin{enumerate}
  \item \textbf{Transparenz.} Die Parlamente, um die es hier geht, sind Parlamente von großen
    Einheiten, wie ganzen Ländern. Bei solchen Parlamenten ist es eigentlich nicht vorstellbar,
    dass eine Einzelperson die nötige Expertise und Bekanntheit erlangt ohne eine Organisation
    im Rücken, die ihr zuarbeitet. Wenn jemand als Einzelbewerber antreten möchte, dann sollte
    seine Organisation auch transparent gemacht werden, indem der Kandidat eine eigene Partei
    gründet.

  \item \textbf{Verfahrensgründe.} Sowohl das oben vorgestellte iterative Wahlverfahren als auch
    das Wahlkreisrepräsentationsverfahren, welches später in diesem Dokument eingeführt wird,
    arbeiten auf Parteienebene. Je nach Stimmenergebnis erhält eine Partei mehr oder weniger
    Sitze im Parlament und vertritt mehr oder weniger Wahlkreise. Wie soll ein einzelner Kandidat
    mehr oder weniger Sitze im Parlament erhalten? Einzelkandidaten würden das ganze Verfahren
    durcheinander bringen.
\end{enumerate}

\quasihead{Schlussfolgerung:}

Eine Auswahl von Personen durch den Wähler muss innerhalb von Parteien vorgenommen werden. Dem
Wähler muss eine Möglichkeit geboten werden, zwischen mehreren Kandidaten derselben Partei eine
Auswahl zu treffen.

\quasihead{Anmerkung:}

Aus dem Verbot von Einzelkandidaten folgt nicht, dass es keine unabhängigen Abgeordneten geben
könnte. Wenn es während einer Legislaturperiode zu einem Zerwürfnis zwischen einem Abgeordneten
und der ihn entsendenden Partei kommt, darf dieser Abgeordnete sein Mandat weiterhin als nunmehr
unabhängiger Abgeordneter ausüben. Diese Erlaubnis beeinträchtigt weder die Transparenz noch das
Verfahren, ist aber unbedingt notwendig, um sicherzustellen, dass die Loyalität eines Abgeordneten
zuerst seinem Gewissen und dem Wähler gelten kann und erst an zweiter Stelle Parteiinteressen
folgen. Wenn diese Reihenfolge umgekehrt würde, würde die parlamentarische Demokratie zur
Diktatur entarten.

\subsection{Wahl der Spitzenkandidaten durch das Volk}

Gerade bei den großen, alteingesessenen Parteien gibt es das Phänomen, dass innerhalb der
Parteien Menschen Karriere machen, die vom Wähler nicht gemocht werden. Das ist kein Zufall. Um
in einer Partei Karriere zu machen, muss man schließlich einer anderen Gruppe Menschen gefallen
bzw.\ sich dort durchsetzen als bei einer Parlamentswahl. Dadurch steht man als Wähler oft vor
der Entscheidung, ob man eine Partei trotz oder wegen ihres Spitzenkandidaten wählen soll. Oftmals
stimmt man inhaltlich mit einer Partei überein, kann sich aber nicht mit deren
Spitzenkandidaten identifizieren. Die CDU hat die Bundestagswahl 2021 deswegen verloren. Umgekehrt
kann es natürlich auch sein. Man mag den Kandidaten, kann aber die Positionen seiner Partei nicht
mittragen.

Zur Abmilderung des Problems schlage ich vor, dass alle sich zur Wahl stellenden Parteien mit
drei Spitzenkandidaten antreten müssen. Jeder Wähler erhält eine \enquote{Parteienstimme}, mit
der er über die Mehrheitsverhältnisse im Parlament bestimmt und eine
\enquote{Spitzenkandidatenstimme}, mit der er unter den Spitzenkandidaten seiner Partei auswählen
kann. Auf dem Stimmzettel steht eine Tabelle mit 4 Spalten: In der ersten Spalte stehen die
Parteien. Daneben stehen jeweils die drei Spitzenkandidaten einer Partei. Der Wähler macht sein
Kreuz bei einer Partei und kann auch ein Kreuz bei einem der Kandidaten der von ihm gewählten
Partei machen. Ein Kreuz bei einem Kandidaten einer anderen Partei ist nicht zulässig. Es gewinnt
(ggf.\ nach mehreren Iterationen) eine Partei aufgrund der Parteienstimmen und der Spitzenkandidat
dieser Partei, der die meisten Stimmen auf sich vereinen konnte.

\subsection{Verfahren zur Direktwahl von Abgeordneten ohne die Parteienzusammensetzung zu
  verändern}

Im aktuellen deutschen Wahlrecht ist eine Beeinflussung der personellen Zusammensetzung des
Bundestages durch den Wähler vorgesehen. Mit seiner \enquote{Erststimme} konnte man früher einen
Kandidaten per Direktmandat unabhängig vom Zweitstimmenergebnis in den Bundestag befördern. Wenn
Erst- und Zweitstimmenergebnis sich dann unterschieden, entstanden Ausgleichsmandate, um den
Proporz zu wahren, die zu einer Vergrößerung des Bundestages führten. Nach dem neuen Wahlrecht
zählen Erststimmengewinne nur, soweit sie vom Zweitstimmenergebnis gedeckt sind. Per Zweitstimme
zustehende Mandate werden zuerst mit Wahlkreisgewinnern (sofern vorhanden) in der Reihenfolge
ihrer relativen Stimmanteile besetzt, bevor mit den Landeslisten weitergemacht wird. Dadurch kann
der Wähler quasi Kandidaten vorne auf die Landesliste setzen und damit theoretisch die personelle
Zusammensetzung der Bundestagsfraktionen beeinflussen. Praktisch ergeben sich dadurch aber nur
Änderungen, wenn die Parteien ihre Wahlkreiskandidaten nicht sowieso vorne auf die Listen setzen.

An dem alten Verfahren wurde vor allem kritisiert, dass der Bundestag durch Überhang- und
Ausgleichsmandate aufgebläht wurde, wenn die Wähler von ihrem Recht des Stimmensplittings
Gebrauch gemacht haben. Dieses Problem kann bei dem neuen Wahlrecht nicht mehr auftreten. Der
Bundestag hat immer fest 630 Sitze.

An dem neuen Wahlrecht wurde kritisiert, dass es Wahlkreise geben kann, die nicht im Bundestag
repräsentiert sind, da sie von Parteikandidaten gewonnen wurden, deren Parteien nicht genügend
Zweitstimmenanteile hatten. Das wäre nach dem alten Wahlrecht nicht passiert, da ein
Wahlkreisgewinn einen sicheren Platz im Bundestag bedeutete.

Um zu sehen, ob man hier noch etwas verbessern kann, muss man zunächst einmal einen Schritt
zurück machen und sich überlegen, was man mit diesem \enquote{Abgeordnetenwahlrecht} überhaupt
erreichen will und kann.

\begin{enumerate}
  \item Man möchte, dass im Parlament Abgeordnete sitzen, die von Wählern persönlich gekannt
    werden und ihnen rechenschaftspflichtig sind. Das erhöht sowohl die Legitimation als auch die
    Kontrolle des Parlaments durch die Bürger. Wenn man in seinem Wahlkreis noch einmal
    wiedergewählt werden will, kann man ihnen nicht das Blaue vom Himmel versprechen und hinterher
    das Gegenteil tun. Mindestens in den ländlichen Gebieten kennen die Wähler ihre Abgeordneten
    auch, und werden nicht aus Versehen Verbrecher ins Parlament wählen. Wenn sie Verbrecher ins
    Parlament wählen, tun sie das mit Absicht, weil sie ihre Gesinnung teilen.

  \item Man möchte alle Gebiete auf dem Hoheitsgebiet des Parlaments auch im Parlament vertreten
    wissen.
\end{enumerate}

Den ersten Punkt kann man verschiedenartig streng auffassen. Man kann a) verlangen, dass im
Parlament genügend vom Bürger persönlich gekannte Abgeordnete sitzen, so dass eine Kontrolle des
Parlaments gegeben ist. Oder man verlangt b), dass alle Abgeordnete des Parlaments vom Bürger
persönlich gekannt werden. Ich halte Forderung~b) für zu stark, da im Parlament ja auch
Sachkompetenz vorhanden sein soll. Während der Bürger Vertrauenswürdigkeit sehr wohl einschätzen
kann und muss, ist das bei Sachkompetenz im Allgemeinen nicht gegeben. Das müssen die Parteien
selbst tun, die ihre Experten ja genau kennen. Ein Experte muss auf seinem Fachgebiet gut sein,
aber nicht notwendigerweise in Bürgerkommunikation. Deswegen würde ich vorschlagen, sich mit
Forderung~1a) zu begnügen und nur die Hälfte des Parlaments mit vom Bürger persönlich gekannten
Kandidaten zu besetzen. (So ist es ja auch ungefähr im aktuellen Wahlrecht.)

Der zweite Punkt lässt sich bei landesweitem Proporz und konstanter Größe des Parlaments im
Allgemeinen nur erfüllen, wenn man unter \enquote{Vertretung der Gebiete} nicht
\enquote{Vertretung der lokalen Macht} versteht.
Das landesweite\footnote{\enquote{Landesweit} oder \enquote{Land} bezieht sich im Folgenden auf
das Hoheitsgebiet des jeweiligen Parlaments. Bei einer Bundestagswahl wäre diese das
Bundesgebiet.} Kräfteverhältnis der Parteien entspricht in der Regel nicht der lokalen Verteilung
der Parteien. Wenn z.\,B.\ Partei~A 60\,\% der landesweiten Wählerstimmen auf sich vereint,
Partei~B 30\,\% und Partei~C 10\,\%, heißt das normalerweise nicht, dass in 60\,\% der Wahlkreise
Partei~A die stärkste Kraft ist, in 30\,\% der Wahlkreise Partei~B und in 10\,\% der Wahlkreise
Partei~C. Wenn alle Parteien landesweit antreten und wirklich die Interessen des ganzen Landes
vertreten, sollte man erwarten, dass Partei~A in jedem Wahlkreis die stärkste Kraft wird. Genau
das ist ja die Stärke des Proporz, da auch bei solch starker landesweiter Überlegenheit einer
Partei die anderen noch im Parlament vertreten sind. Die Erfahrung zeigt, dass auch landesweit
antretende Parteien oft regionale Hochburgen haben, in denen sie selbst bei Wahlen traditionell
stark abschneiden und ihre politischen Gegner entsprechend schwächer. Deswegen könnte es sein,
dass doch auch beim obigen Beispiel Partei~B oder~C noch den einen oder anderen Wahlkreis holen.
Bedenklich für den Zusammenhalt eines Landes wäre es allerdings, wenn tatsächlich 60\,\% aller
Wahlkreise Partei~A wählen, 30\,\% aller Wahlkreise Partei~B und 10\,\% alle Wahlkreise Partei~C
und das auch noch in zusammenhängenden Gebieten. Das würde ja bedeuten, dass offensichtlich keine
Partei die Interessen des ganzen Landes im Blick hat, sondern nur ihr jeweiliges Kerngebiet.
Darüber hinaus würde Gebiet A quasi über die Gebiete B und C herrschen. Das liefert Zündstoff für
Abspaltungsbestrebungen oder gar Bürgerkriege.
In diesem Sinne ist das Ergebnis der Bundestagswahl 2025 sehr bedenklich, wo der ganze Osten blau
gewählt hat, und der Westen schwarz mit einigen roten Enklaven.\footnote{\url{https://www.tagesschau.de/inland/bundestagswahl/wahlkreisergebnisse-karte-bundestagswahl-2025-100.html}} Das ist eine
Tendenz, die sich nicht fortsetzen darf!
Wenn wir aber von einer \enquote{gesunden} regionalen Stimmenverteilung ausgehen, d.\,h.\ die
Kräfteverhältnisse zwischen den Parteien sind im ganzen Land etwa gleich abgesehen von ein paar
regionalen Schwankungen (\enquote{Hochburgen} oder Zufällen), müssen wir auch Parlamentssitze von
Parteien mit Kandidaten befüllen, die in keinem Wahlkreis gewonnen haben und dafür
\enquote{rechtmäßigen Wahlkreisgewinnern} ihren Wahlkreis wegnehmen. Dieser Fehler ist
prinzipbedingt. Er lässt sich nicht vermeiden, höchstens minimieren. Irgendwo müssen wir den
Fehler hinschieben: Entweder wir machen die Größe des Parlaments variabel, oder wir leben damit,
dass einige Wahlkreise keine Repräsentation im Parlament haben, oder wir nehmen einigen
Wahlkreisgewinnern ihren Wahlkreis weg und geben sie einem Verlierer, damit das Sitzverhältnis im
Parlament dem landesweiten Wählerwillen entspricht. Parlamentsvergrößerungen sind schon seit
langem unpopulär, fehlende Representation einiger Wahlkreise wurde bei der Bundestagswahl 2025
kritisiert. Es bleibt nur übrig, einigen Wahlkreisgewinnern ihren Wahlkreis wieder wegzunehmen.
Eine Abkehr vom landesweiten Proporz wäre auch noch eine Möglichkeit. Wenn die
Parteienzusammensetzung des Parlaments durch Parteizugehörigkeit der Wahlkreisgewinner bestimmt
würde, würde jeder Wahlkreis repräsentiert, aber die Mehrheitsverhältnisse im Parlament würden
nicht mehr das landesweite Kräfteverhältnis der Parteien repräsentieren. Ein iteratives
Verfahren, wie wir es auch vorgeschlagen haben, würde dann auch keinen Sinn mehr machen, da es
ja die Mehrheitsverhältnisse wieder verändern würde.

Der Ausweg besteht darin, nur zu verlangen, dass jedes Gebiet eine Vertretung bei der
landesweiten Macht hat und nicht, dass die lokalen Kräfteverhältnisse im Parlament
widergespiegelt werden müssen.

\enquote{Vertretung bei der landesweiten Macht} kann man verschiedenartig streng auffassen:
\begin{enumerate}[label=\alph*)]
  \item Jeder Wahlkreis hat einen Vertreter in der Regierungspartei bzw.\ der
    Regierungskoalition.
  \item Jeder Wahlkreis entsendet einen Vertreter ins Parlament. Es muss sich aber nicht um einen
    Vertreter in der Regierungskoalition handeln. Einige Wahlkreise müssen sich mit einer
    parlamentarischen Vertretung durch Oppositionspolitiker begnügen.
\end{enumerate}

Forderung~a) ist stärker und gerechter als Forderung~b), da ein Zugang zur
Regierungspartei/koalition natürlich potentiell mehr Einflussmöglichkeiten mit sich bringt und
Wahlkreise ohne Regierungsvertretung, die es bei~b) ja geben würde, gegenüber denen mit
Regierungsvertretung benachteiligt wären. Forderung~a) erfordert aber auch ungefähr doppelt so
große Parlamente wie Forderung~b).
Wir hatten ja weiter oben aus gutem Grund verlangt, dass nur die Hälfte des Parlaments direkt
vom Bürger gewählt wird. Die andere Hälfte sollte von den Parteien besetzt werden, damit diese
ihre Experten einbringen können, die nicht unbedingt gut in Bürgerkommunikation sein müssen. Wenn
Forderung~a) zur Anwendung kommt, dann muss dieses Prinzip auch für die Regierungspartei/koalition
gelten, die i.\,A.\ nur leicht über die Hälfte der Parlamentssitze verfügt. Wenn garantiert immer
alle Wahlkreise eine Regierungsvertretung haben sollen, dann darf die Anzahl der Wahlkreise nicht
mehr als ein Viertel der Parlamentssitze betragen. Die Wahlkreise sollten aber so eine Größe
besitzen, dass es noch möglich ist, Abgeordnete persönlich zu kennen. Nach den Erfahrungen des
Autors ist dies bei 200\,000 -- 300\,000 Einwohnern pro Wahlkreis tatsächlich auch noch der Fall.
Aus der maximalen Anzahl der wahlberechtigten Einwohner pro Wahlkreis und der Anzahl der
Wahlberechtigten insgesamt ergibt sich aber die minimale Anzahl der Wahlkreise, woraus sich
wieder die minimale Parlamentsgröße zu a) $4\times$ minimale Anzahl der Wahlkreise und in
Variante~b) zu $2\times$ minimale Anzahl der Wahlkreise ergibt.

In Deutschland haben wir zur Zeit 299 Wahlkreise, die im Mittel ca.\ 240\,000 Einwohner stark
sind.\footnote{\url{https://www.bundestag.de/parlament/bundestagswahl/wahlkreise-1050626}} Bei
630 Parlamentssitzen würde also Variante~a) nicht umsetzbar sein. Nun könnte man sowohl die
Einwohnerzahl pro Wahlkreis als auch die Parlamentsgröße erhöhen, um Variante~a) umsetzen zu
können, oder man nimmt die durch Variante~b) entstehende Ungleichbehandlung der Wahlkreise
zugunsten eines kleineren Parlaments in Kauf.

Die Regierung eines Staatsgebiets muss in allen drei Gewalten unbedingt im Interesse des gesamten
Landes handeln und ggf.\ Abwägungen zwischen sich widersprechenden regionalen Interessen treffen.
Auch wenn durch die Wahlkreisabgeordneten doch durch die Hintertür wieder eine Art
Regionalvertretung ins Spiel kommt, so sollte diese doch eine untergeordnete Rolle spielen, da ja
in einem Landesparlament\footnote{nicht \enquote{Bundeslandesparlament}}, wie dem Deutschen
Bundestag, ausdrücklich Politik gemacht und Gesetze beschlossen werden, die für das ganze Land
gelten. Weiterhin gab es auch bei dem bisherigen System der Direktmandate keinerlei Anspruch auf
eine Gebietsvertretung durch einen Regierungspartei/koalitionsvertreter. Vielmehr konnte man
seinen Abgeordneten ins Parlament wählen. Ob der dann zur Regierungsmehrheit gehörte oder nicht,
hing von den Kräfteverhältnissen im Bundestag ab. Niemand hat die daraus entstehende
Gebietsungleichbehandlung jemals gestört. Deshalb denke ich, man könnte mit Variante~b) leben.
Am Ende ist es aber eine gesamtgesellschaftliche Entscheidung, wie viel uns eine Gleichbehandlung
von Gebieten wert ist. Solche Fragen könnte man gegebenenfalls auch dem Volk zur Abstimmung
vorlegen.

Unabhängig davon, für welche Variante man sich entscheidet, bleibt das Problem, den lokalen
Wahlkreisvertreter zu bestimmen, ohne die Parteienzusammensetzung im Parlament zu verändern. Der
Vorschlag ist, dass jede Partei in jedem Wahlkreis mit mindestens drei Kandidaten antreten muss.
Jeder Wähler kann bei jeder Partei in seinem Wahlkreis einen Kandidaten als seinen
Wunschabgeordneten bestimmen. Er kann das für jede antretende Partei tun, nicht nur für die, die
er mit seiner \enquote{Parteienstimme} gewählt hat. Jeder Wähler bekommt also 3 Arten von
Stimmen:
\begin{enumerate}
  \item Eine \enquote{Parteienstimme}, mit der er eine Partei ins Parlament wählt.
  \item Eine \enquote{Spitzenkandidatenstimme}, die er für einen Spitzenkandidaten seiner mit~1.
    gewählten Partei abgeben kann.
  \item Eine \enquote{Abgeordnetenstimme} pro antretender Partei. Für jede antretende Partei kann
    er in seinem Wahlkreis seinen Lieblings- oder am wenigsten gehassten Kandidaten wählen.
\end{enumerate}
Nachdem die Parteienzusammensetzung des Parlaments feststeht, wird nach einem mathematischen
Verfahren, das die relativen Mehrheiten der Parteien pro Wahlkreis berücksichtigt und bei
Nicht-Eindeutigkeit durch ein Zufallsverfahren ergänzt wird, jeder Wahlkreis genau einer Partei
zugeteilt. Der Kandidat dieser Partei, der die meisten Abgeordnetenstimmen auf sich vereint, ist
der Abgeordnete dieses Wahlkreises. Es ist also nicht möglich, über die Wahl einer Person, das
zahlenmäßige Verhältnis der Parteien im Parlament noch einmal zu ändern, nachdem es durch die
Parteienstimmen bereits festgelegt wurde. Trotzdem ist es möglich, die personelle Zusammensetzung
des Parlaments zu beeinflussen, indem man selbst bei fremden Parteien von seiner
Abgeordnetenstimme Gebrauch macht.

\subsection{Verfahren zur Zuweisung eines Wahlkreises zu einer Partei
  (Wahlkreisrepräsentationsverfahren)}

Ein Wahlgebiet $W$ (z.\,B.\ ein Land) sei in $M$ Wahlkreise $w_1 \ldots w_M$ aufgeteilt. Jedem
Wahlkreis entspricht ein vom Bürger nach obigen Verfahren per Direktwahl zu besetzender
Parlamentssitz. Es sind also $M$ Parlamentssitze zu verteilen. Da wir ausdrücklich auch
Koalitionen zur Regierungsbildung zulassen, nur nicht mehr erzwingen wollen, ist es für das
Verfahren egal, ob $M$ der Hälfte des Parlaments oder nur der Hälfte der Regierungskoalition
entspricht. Das Verfahren muss $M$ Wahlkreise auf $N$ Parteien aufteilen können. Im Falle, dass
obige Variante~a) angewendet wird und eine Partei die absolute Mehrheit erringt oder ihr
zugesprochen wird, ist $N$ eben 1.

Aus dem Ergebnis der Parteistimmen ergibt sich eine Direktmandatsverteilung $\mu$, die den
Parteien $p_1 \ldots p_N$ die Anzahl Mandate (Parlamentssitze) $m_1 \ldots m_N$ zuordnet.
Welches Verfahren verwendet wird, um aus dem Parteistimmenverhältnis bei gegebener Gesamtzahl $M$
die Direktmandatsverteilung $\mu$ zu bestimmen, spielt für die folgenden Ausführungen keine Rolle,
wie man ein rationales Stimmenverhältnis durch eine ganzzahlige Sitzverteilung am besten
approximieren kann. Im Laufe der Zeit sind verschiedene Verfahren für dieses Problem entwickelt
worden, deren Vor- und Nachteile schon so lange Gegenstand der Forschung sind. Als Einstiegspunkt
in die Materie mag z.\,B.\ \enquote{Wahlrecht und Parteiensystem} von Dieter Nohlen herangezogen
werden.

Im Folgenden wird $m_1 \ldots m_N$ als gegeben vorausgesetzt.

Die Parteien $p_1 \ldots p_N$ erreichen in den Wahlkreisen $w_1 \ldots w_M$ jeweils die Stimmen

\[
  \underline{S} = \begin{pmatrix}
    s_{11} & \cdots & s_{1N} \\
    \vdots & s_{ij} & \vdots \\
    s_{M1} & \cdots & s_{MN}
  \end{pmatrix}
\]

mit $s_{ij}$ -- Anzahl der Parteienstimmen im $i$-ten Wahlkreis für die $j$-te Partei.
Interessant sind aber eigentlich nur die relativen Stimmanteile pro Wahlkreis

\[
  \underline{R} = \begin{pmatrix}
    r_{11} & \cdots & r_{1N} \\
    \vdots & r_{ij} & \vdots \\
    r_{M1} & \cdots & r_{MN}
  \end{pmatrix}
\]

mit $r_{ij} = s_{ij}/S_i$, wobei $S_i$ für die Gesamtzahl der abgegebenen Parteienstimmen im
$i$-ten Wahlkreis steht. Die Summe der relativen Stimmanteile in einem Wahlkreis ist natürlich
eins. Ausgedrückt in Formeln:
\[
  S_i = \sum_{j=1}^{N} s_{ij}
  \qquad \text{und} \qquad
  \sum_{j=1}^{N} r_{ij} = \sum_{j=1}^{N} \frac{s_{ij}}{S_i}
    = \frac{1}{S_i}\sum_{j=1}^{N} s_{ij} = 1
\]

Gesucht ist ein Verfahren $V(R)$, das jedem Wahlkreis $w_i$ eindeutig eine Partei $p_j$ zuordnet.
Es soll $N \leq M$ gelten, d.\,h.\ es kann passieren, dass eine Partei mehreren Wahlkreisen
zugeordnet wird, aber es kann nie vorkommen, dass eine Partei gar keinen Wahlkreis abbekommt,
weil zu wenig Wahlkreise zur Verfügung stehen.\footnote{Diese Bedingung wird in der Praxis von
Bundestagswahlen wohl immer erfüllt sein, ist aber streng genommen nicht notwendig. Wenn $N > M$
gilt, kann man immer die Parteien absteigend nach ihren Parteienstimmenergebnissen ordnen und die
Wahlkreise auf die ersten $N' = M$ Parteien aufteilen. Die übrigen Parteien sind dann zwar im
Parlament vertreten, aber repräsentieren keinen Wahlkreis mehr.}

Für den Fall $N = 1$ ist das Verfahren trivial. Alle Wahlkreise werden der einzig zur Verfügung
stehenden Partei zugeordnet.

Für $N = 2$ kann man die Wahlkreise aus Sicht jeder Partei absteigend nach ihren relativen
Stimmanteilen sortieren. Also könnte sich z.\,B.\ für $M = 5$ für Partei $p_1$ folgende Ordnung
$\Omega_1$ ergeben:
\[
  \Omega_1 = [w_3, w_5, w_2, w_4, w_1]
\]
Das würde bedeuten, dass $r_{31} > r_{51} > r_{21} > r_{41} > r_{11}$ gilt.

Genauso könnte man jetzt $\Omega_2$ als nach relativem Wahlergebnis geordnete Liste für $p_2$
bestimmen. Für $N = 2$ ist $\Omega_2$ aber schon durch implizit durch $\Omega_1$ als $\Omega_1$ in
umgekehrter Reihenfolge gegeben, da $r_{i1} + r_{i2} = 1$ gilt. Also $r_{i2} = 1 - r_{i1}$. Da,
wo die erste Partei relativ die meisten Stimmen geholt hat, hat die zweite Partei die wenigsten
und umgekehrt.

Also für das obige Beispiel würde sich $\Omega_2$ zu $\Omega_2 = [w_1, w_4, w_2, w_5, w_3]$
ergeben.

Diese $\Omega$-Listen sagen aus, wie wichtig diese Partei den Bürgern in einem Wahlkreis ist, im
Verhältnis zu der Wichtigkeit, die ihr Bürger in anderen Wahlkreisen zumessen. Daraus lässt sich
natürlich auch umgekehrt eine Wichtigkeit der Wahlkreise für die Partei ableiten: Wenn eine Partei
nur einen Teil der Wahlkreise vertreten kann (was für $N > 1$ immer der Fall ist), dann sollte sie
natürlich möglichst den Teil vertreten, wo sie mehr gewollt wird, als in dem anderen Teil. Für
$N = 2$ lässt sich das auch immer erfüllen, da die Wichtigkeit ja genau umgekehrt ist. Für
$N = 2$ kann man jeder Partei einfach die ersten $m_j$ Wahlkreise aus der $\Omega_j$-Liste
zuteilen.

Dann werden alle Wahlkreisen zugeteilt sein, und jede Partei vertritt den für sie wichtigsten
Teil der Wahlkreise.

Was machen wir im Fall $N > 2$? Für $N > 2$ kann ja derselbe Wahlkreis durchaus für mehrere
Parteien der wichtigste sein, nur nicht für alle gleichzeitig. Man kann sich das anschaulich gut
vorstellen. Im unteren Bild ist eine beispielhafte Verteilung der relativen Stimmanteile für drei
Parteien dargestellt, wobei die 2.\ und 3.\ Partei beide gleichzeitig ihre maximalen relativen
Stimmanteile da haben, wo die erste Partei ihr Minimum hat. (Auf der $y$-Achse sind die Wahlkreise
aufgetragen. Die Abstände zwischen den Linien bzw.\ den Rändern sind die relativen Stimmanteile
der Parteien. Die Summe der relativen Stimmanteile ist 1. Der Einfachheit halber nehmen wir an,
es seien so viele Wahlkreise, dass man die Lücken zwischen den Punkten nicht mehr sieht und wir
somit die Linien durchzeichnen können.)

\bigskip
\begin{center}
\includegraphics[width=0.8\textwidth]{diagrams/maximum_for_2_and_3_same_constituency.pdf}
\end{center}
\bigskip

Es hätte auch anders sein können. Im folgenden Beispiel haben alle 3 Parteien ihre maximalen
relativen Stimmanteile in verschiedenen Wahlkreisen.

\bigskip
\begin{center}
\includegraphics[width=0.8\textwidth]{diagrams/maxima_in_different_constituencies.pdf}
\end{center}
\bigskip

Da wir für $N = 2$ bereits ein funktionierendes Verfahren haben, könnte man versuchen, das
Verfahren für $N > 2$ auf $N = 2$ zurückzuführen, indem man die Parteien nach ihren
Parteienstimmen absteigend ordnet und alle Parteien außer der ersten zu einer Partei $p'$
zusammenfasst. Dann könnte man der ersten Partei nach obigem Verfahren ihre Wahlkreise zuordnen.
Im nächsten Schritt spaltet man $p'$ in $p_2$ und $p''$ auf und wiederholt das Verfahren mit den
noch nicht zugewiesenen Wahlkreisen. Das macht man immer wieder, solange bis alle Wahlkreise
zugewiesen sind. Dieses Verfahren hat nur den Nachteil, dass die jeweils parteienstimmenstärkste
Partei den anderen Parteien ihre relative beste Wahlkreise wegnehmen kann. Im folgenden Beispiel
hat die kleine Partei $p_3$ ihr bestes Ergebnis ausgerechnet da, wo auch die große Partei $p_1$
ihr bestes Ergebnis hat. Aber $p_1$ hat im Gegensatz zu $p_3$ eine \enquote{dicken Bauch}. Selbst
wenn sie ihre stärksten Wahlkreise an $p_3$ abgeben würde, ist sie in den anderen Wahlkreisen
auch immer noch stark vertreten, was bei $p_3$ nicht der Fall ist. Da gehen die relativen
Stimmanteile mit steigendem $i$ doch ziemlich schnell runter. Deshalb sollte $p_3$ die ersten
Wahlkreise bekommen und nicht $p_1$, obwohl sie dort stärkste Kraft ist.

\bigskip
\begin{center}
\includegraphics[width=0.8\textwidth]{diagrams/thick_belly.pdf}
\end{center}
\bigskip

Wir brauchen ein Maß für die \enquote{relative Wichtigkeit} eines Wahlkreises für eine Partei,
das sich parteiübergreifend vergleichen lässt. Dazu können wir aus der Matrix $R$ die Matrix
$\Psi(R)$ bestimmen, für deren Elemente $\psi_{ij}$ gelten soll:

\[
  \psi_{ij} = (M-1)\,\frac{r_{ij}}{\displaystyle\sum_{\substack{i' \neq j,\, i'=1}}^{M} r_{i'j}}
  \quad \text{für} \quad \sum_{\substack{i \neq j,\, i=1}}^{M} r_{ij} \neq 0
  \quad \text{und $K$ sonst,}
\]

wobei $K$ eine beliebig gewählte sehr große Zahl sein soll ($K \gg \psi_{ij}\ \forall\,i,j$).

$\psi_{ij}$ gibt das Verhältnis des relativen Stimmenanteils im Wahlkreis $i$ für die Partei~$j$
zur Summe ihrer relativen Stimmenanteile in den anderen Wahlkreisen an. Theoretisch könnte diese
Summe auch 0 sein, wenn die Partei~$j$ tatsächlich in allen Wahlkreisen außer dem $i$-ten keine
Stimmen erhalten hat. Da jede Partei in jedem Wahlkreis antreten muss, ist das sehr
unwahrscheinlich. Für diesen unwahrscheinlichen Fall kann man dann einfach eine sehr große Zahl
$K$ festlegen, die man dem Element anstelle von unendlich zuweist. Je wichtiger ein Wahlkreis für
eine Partei im Vergleich zu allen anderen Wahlkreisen ist, desto größer wird $\psi_{ij}$. Der
Faktor $(M-1)$ stellt lediglich eine Normierung dar, damit die Zahlenwerte von $\Psi$ besser
menschenlesbar sind. Durch diese Normierung ist $\Psi(R)$ bei Gleichverteilung überall 1.

Um nun jeden Wahlkreis einer Partei zuzuweisen sucht man zunächst aus $\Psi$ das größte Element
$\psi_{ij}$. Wenn es mehrere gleich große $\psi_{ij}$ gibt, wählt man eines davon zufällig aus.
Den zugehörigen Wahlkreis $w_i$ weist man der Partei $p_j$ zu. Die Anzahl der dieser Partei noch
zuzuweisenden Mandate $m_j$ verringert man um eins und die Zeile $w_i$ wird als nicht mehr
auswählbar markiert.

Jetzt wiederholt man das Verfahren mit den noch auswählbaren $\psi_{ij}$ und den Parteien, die
noch zuzuweisende Mandate $m_j > 0$ haben. Spalten von Parteien mit $m_j = 0$ werden als nicht
auswählbar markiert. Das Verfahren ist beendet, wenn $m_j = 0$ für alle Parteien gilt. D.\,h.\
keine Spalte ist mehr auswählbar.

Es werden also pro Partei nur die durch das Parteienstimmenergebnis bestimmte Anzahl Mandate
verteilt. Dabei wird mit der Verteilung bei den regional am ungleichmäßigsten verteilten Parteien
begonnen, weil diese bestimmte Wahlkreise am dringendsten brauchen.
Gleichmäßiger verteilte Parteien sind nicht so sehr auf einen bestimmten Wahlkreis angewiesen.
Sie haben überall gleich gut bzw.\ gleich schlecht abgeschnitten.

\quasiheadul{Zahlenbeispiel:}

Es treten 3 Parteien $p_1 \ldots p_3$ in 5 Wahlkreisen $w_1 \ldots w_5$ an. Der Einfachheit
halber sollen alle Wahlkreise dieselbe Größe $G = 1000$ haben. Die folgende Tabelle listet die
für die jeweiligen Parteien im jeweiligen Wahlkreis abgegebenen Parteienstimmen auf:

\bigskip
\begin{center}
\begin{tabular}{l r r r}
  \toprule
           & $p_1$ & $p_2$ & $p_3$ \\
  \midrule
  $w_1$    & 650   & 100   & 250   \\
  $w_2$    & 630   & 200   & 170   \\
  $w_3$    & 600   & 280   & 120   \\
  $w_4$    & 550   & 350   & 100   \\
  $w_5$    & 450   & 500   & 50    \\
  \midrule
  $\Sigma$ & 2880  & 1430  & 690   \\
  \bottomrule
\end{tabular}
\end{center}
\bigskip

Es sind fünf Wahlkreise vorhanden, also müssen 5 Mandate irgendwie anhand der Wahlergebnisse auf
3 Parteien aufgeteilt werden. Vom Prinzip her muss man die relativen Stimmanteile $P_j$ der
Parteien auf das ganze Wahlgebiet gerechnet bestimmen, mit der Anzahl der Mandate multiplizieren
und dann geeignet runden. Also:
\[
  P_j = \frac{1}{S}\sum_{i=1}^{M} s_{ij}
  \quad \text{mit} \quad S = \sum_{i=1}^{M}\sum_{j=1}^{N} s_{ij}
\]
\[
  m_j = \textit{runde\_geeignet}(P_j \cdot M)
\]
Aber was bedeutet \enquote{geeignet runden}? Bei einfacher Standardrundung ergibt die Summe der
$m_j$ nicht unbedingt $M$. Es können durch das Runden mehr oder weniger Mandate herauskommen, als
zu verteilen sind.

Deswegen sind verschiedene Verfahren entwickelt worden, um aus einer Stimmenverteilung auf eine
Mandatsverteilung zu kommen. Bei der Bundestagswahl kommt das Verfahren nach Sainte-Lagüe/Schepers
zum Einsatz.\footnote{\url{https://de.wikipedia.org/wiki/Sainte-Lagu\%C3\%AB-Verfahren},
\url{https://www.bundeswahlleiterin.de/service/glossar/s/sainte-laque-schepers.html}}

Das Verfahren nach Sainte-Lagüe/Schepers ergibt für das obige Beispiel:
$m_1 = 3$, $m_2 = 1$ und $m_3 = 1$.

Also 3 Wahlkreise gehen an Partei $p_1$ und jeweils ein Wahlkreis an Partei $p_2$ und $p_3$. Um
zu bestimmen, welche Wahlkreise an welche Partei gehen, stellen wir zunächst $\underline{R}$ auf.
Da in unserem Beispiel die Wahlkreisgröße konstant 1000 sein soll, ist das einfach der Inhalt
der obigen Tabelle geteilt durch 1000:

\[
  \underline{R} = \begin{pmatrix}
    0.65 & 0.10 & 0.25 \\
    0.63 & 0.20 & 0.17 \\
    0.60 & 0.28 & 0.12 \\
    0.55 & 0.35 & 0.10 \\
    0.45 & 0.50 & 0.05
  \end{pmatrix}
\]

Daraus bestimmen wir $\Psi$ zu

\[
  \underline{\Psi} = \begin{pmatrix}
    1.1659 & 0.3008 & 2.2727 \\
    1.1200 & 0.6504 & 1.3077 \\
    1.0526 & 0.9739 & 0.8421 \\
    0.9442 & 1.2963 & 0.6780 \\
    0.7407 & 2.1505 & 0.3125
  \end{pmatrix}
\]

Der höchste Wert von $\underline{\Psi}$ ist $\psi_{13} = 2.2727$. D.\,h.\ der Wahlkreis $w_1$
geht an Partei $p_3$, sofern dieser Partei noch Stimmen zustehen, d.\,h.\ sofern $m_3 > 0$. Das
ist es, so dass $w_1 \longrightarrow p_3$, $m_3$ reduziert sich um eins, also $m'_3 = 0$.

D.\,h.\ Partei $p_3$ kriegt keinen anderen Wahlkreis mehr zugewiesen und Zeile 1 kann aus
$\underline{\Psi}$ nicht mehr ausgewählt werden.

Die nächstgrößte Element in $\underline{\Psi}$ ist $\psi_{52} = 2.1505$. Deswegen wird Wahlkreis
$w_5$ Partei $p_2$ zugewiesen und $m_2$ um eins erniedrigt. Wir haben bis jetzt also:
$w_1 \longrightarrow p_3$, $w_5 \longrightarrow p_2$, $m''_1 = 3$, $m''_2 = 0$, $m''_3 = 0$.

In der Matrix sind jetzt nur noch Elemente aus der ersten Spalte auswählbar, die alle an $p_1$
gehen, die bislang ja noch gar keinen Wahlkreis abbekommen hatte. $p_1$ hat auch noch genügend
Mandate, da wir mit diesem Verfahren ja immer genauso viele Mandate wie Wahlkreise verteilen.

\quasiheadul{Endergebnis:}

\begin{itemize}[label={}]
  \item $w_1 \longrightarrow p_3$
  \item $w_2 \longrightarrow p_1$
  \item $w_3 \longrightarrow p_1$
  \item $w_4 \longrightarrow p_1$
  \item $w_5 \longrightarrow p_2$
\end{itemize}

\subsection{Ungarischer Algorithmus zur Optimierung des Zuweisungsproblems}

Der oben vorgeschlagene Algorithmus zur Zuweisung von Parteien zu Wahlkreisen hat ein Problem: Er
ist \enquote{gierig}. Er nimmt sich immer das nächstgrößte Element und übersieht dadurch manchmal
günstigere Kombinationen.

Konstruiertes Beispiel mit unnatürlich hohen Wichtigkeitswerten, aber zeigt das Problem:

Mit $m_1 = 2$, $m_2 = 1$, $m_3 = 1$:

\[
  \underline{\Psi} = \begin{pmatrix}
    5 & 1 & 6 \\
    7 & 5 & 1 \\
    1 & 5 & 4 \\
    1 & 1 & 5
  \end{pmatrix}
\]

Der Gierige Algorithmus würde so vorgehen:
\begin{itemize}[label={}]
  \item $w_2 \longrightarrow p_1$
  \item $w_1 \longrightarrow p_3$
  \item $w_3 \longrightarrow p_2$
  \item $w_4 \longrightarrow p_1$
\end{itemize}

Die Summe der Wichtigkeiten wäre: $\mathrm{WS}_{\mathrm{GREEDY}} = 7 + 6 + 5 + 1 = 19$

Die optimale Lösung (im Sinne von Maximierung der zugewiesenen Gesamtwichtigkeit) wäre aber:
\begin{itemize}[label={}]
  \item $w_1 \longrightarrow p_1$
  \item $w_2 \longrightarrow p_1$
  \item $w_3 \longrightarrow p_2$
  \item $w_4 \longrightarrow p_3$
\end{itemize}

$\mathrm{WS}_{\mathrm{OPTIMAL}} = 5 + 7 + 5 + 5 = 22$

Das Problem ist, dass der Gierige Algorithmus nicht vorausschauend ist. Durch die Zuweisung
$w_1 \to p_3$ (Wert 6 statt 5 für $p_1$) gewinnt er kurzfristig 1 Punkt, verliert aber in Zeile
4 vier Punkte (Wert 1 statt 5 für $p_1$). Deswegen ist die Summe der Wichtigkeiten beim Gierigen
Algorithmus in diesem Beispiel $4 - 1 = 3$ schlechter als die optimale Lösung.

Es gibt bereits seit langem einen Standardalgorithmus für derartige Zuweisungsprobleme: Der
sogenannte \href{https://de.wikipedia.org/wiki/Ungarische\_Methode}{Ungarische Algorithmus}
findet garantiert die optimale Zuweisung im Sinne von Maximierung der Summe der zugewiesenen
relativen Wichtigkeiten. Eine Version dieses Algorithmus kommt auch in der Simulation
standardmäßig zum Einsatz.

\quasihead{Hinweis:} Theoretisch kann es vorkommen, dass es kein globales Maximum gibt, sondern
mehrere gleichwertige Zuweisungen. In diesem Fall liefert der Ungarische Algorithmus irgendeine
davon zurück. Bei realen Wahlen auf Bundes- oder Bundesländerebene ist dieser Fall sehr
unwahrscheinlich, da sich bei vielen abgegebenen Stimmen die Ergebnisse pro Wahlkreis und Partei
immer geringfügig unterscheiden werden.

Der Gierige Algorithmus ist bei optimierter Implementierung nicht so rechenintensiv wie der
Ungarische Algorithmus.\footnote{Algorithmische Ordnung: Gieriger Algorithmus: naive
Implementierung $O(M^2 N)$, optimiert $O(M \cdot N \cdot \log(M \cdot N))$, Ungarischer
Algorithmus: $O(M^2 N)$.} Der Unterschied spielt bei realen Wahlen jedoch praktisch keine Rolle.
Beide Algorithmen liefern auf 08/15-Laptops in Millisekunden Ergebnisse.

Auch wenn der Gierige Algorithmus meistens dieselben Ergebnisse wie der Ungarische Algorithmus
liefert, würde ich, da die Laufzeit praktisch keine Rolle spielt, für so ein sensibles Thema wie
Wahlkreisrepräsentation den Ungarischen Algorithmus empfehlen. Bei ihm ist garantiert, dass eine
optimale Zuordnung gefunden wird.
